\chapter{Open chains with boundary fields}
\label{chap:open}

The open spin-1 Heisenberg chain with equal boundary fields \(h = 3/5\),
\[
  H_n = \sum_{i=0}^{n-2} \mathbf S_i \cdot \mathbf S_{i+1} - h\,(S_0^z + S_{n-1}^z),
\]
has a unique ground state and a spectral gap \(\Delta_n > (\log 10)/392 \approx 0.005874\) for
every odd \(n \ge 1921\) (Theorem~2 of~\cite{HaldaneNote}, Appendices~A--C; the theorem and its
four finite-size inputs are from~\cite{OpenAIBoundary}). This chapter establishes the bound
conditional on those four finite-size inputs.

\section{Transfer-operator boundary moments}

At every finite Trotter number, the ring partition functions and the open-chain partition functions
are spectral moments of a single real symmetric transfer matrix \(X = \mathcal F J\), where
\(\mathcal F\) is the single-cell transfer matrix and \(J\) reverses imaginary time:
\(Z_{2N} = \Tr X^{2N}\) and \(G_n = \langle v, X^n v\rangle\) for a real boundary vector \(v\).
Choosing equal boundary fields at odd chain lengths and opposite boundary fields at even chain
lengths gives both parities this bilinear moment representation with the same boundary vector
\(v\).

\begin{definition}[Transfer-operator boundary moment inequalities]
  \label{def:open-moment-bounds}
  \lean{CertifiedGaps.OpenMomentWitness, CertifiedGaps.OpenMomentBounds,
    CertifiedGaps.leadingRoot, CertifiedGaps.leadingOverlap, CertifiedGaps.boundaryDefect}
  \leanok
  For ring partition functions \(Z\) and open-chain partition functions \(G\), a boundary-moment
  witness \((\mu, c, \rho, s)\) consists of a leading transfer eigenvalue \(\mu \ge 0\), a leading
  boundary overlap \(c \ge 0\), a subleading spectral ratio \(\rho \in [0, 1]\), and a leading
  parity sign \(s = \pm 1\) satisfying the closed moment inequalities induced by a Hermitian
  transfer operator: in particular, \(c\mu^{2N} \le G_{2N}\), the subleading decay bound
  \(\mu^{2K}(G_{2N} - c\mu^{2N}) \le \rho^{2N-2K}\mu^{2N}(G_{2K} - c\mu^{2K})\) for \(K \le N\),
  and the odd-length bound
  \(|G_{2N+1} - sc\mu^{2N+1}| \le \mu^{2N+1-2K}(G_{2K} - c\mu^{2K})\). The \emph{boundary defect}
  is \(e_k = 1 - c\mu^k/G_k\), where \(\mu = \inf_N Z_{2N}^{1/(2N)}\) and
  \(c = \inf_N G_{2N}/\mu^{2N}\).
\end{definition}

\begin{lemma}[Boundary moment inequalities for Hermitian matrices and their limits]
  \label{lem:open-moment-bounds-limit}
  \lean{CertifiedGaps.OpenMomentBounds.of_eigen, CertifiedGaps.OpenMomentBounds.of_hermitian,
    CertifiedGaps.OpenMomentBounds.of_tendsto, CertifiedGaps.OpenMomentWitness.leadingRoot_eq,
    CertifiedGaps.OpenMomentWitness.leadingOverlap_eq}
  \leanok
  \uses{def:open-moment-bounds}
  The moments \(\Tr X^{2N}\) and \(\langle v, X^n v\rangle\) of any Hermitian matrix \(X\) satisfy
  the boundary moment inequalities, and these inequalities are preserved under pointwise limits with
  \(Z_2, Z_4 > 0\). Every witness satisfies \(\mu = \inf_N Z_{2N}^{1/(2N)}\), and
  \(c = \inf_N G_{2N}/\mu^{2N}\) whenever \(\rho < 1\).
\end{lemma}
\begin{proof}
  \leanok
  Take \(\mu\) to be the largest eigenvalue modulus of \(X\), \(c\) its boundary overlap weight,
  and \(\rho\mu\) the next largest eigenvalue modulus; each inequality holds termwise in the
  spectral decomposition. For pointwise limits, the witnesses lie in a compact box determined by
  \(Z_2\), \(Z_4\), and \(G_0\), so a convergent subsequence yields a limiting witness because each
  inequality is closed.
\end{proof}

\begin{theorem}[Boundary moments of reflection-positive open chains]
  \label{thm:open-moment-bounds}
  \lean{CertifiedGaps.RPRing.Boundary, CertifiedGaps.RPRing.openHamiltonian,
    CertifiedGaps.RPRing.openTransfer_isHermitian, CertifiedGaps.RPRing.trace_openSlice_pow,
    CertifiedGaps.RPRing.openMomentBounds}
  \leanok
  \uses{lem:open-moment-bounds-limit, def:rpring, def:cell-transfer, lem:chernoff}
  For a reflection-positive ring with trivial internal reflection (\(F = 1\)) and a boundary
  operator \(b = -\sum_\alpha w_\alpha C_\alpha\) whose coefficients \(w_\alpha\) are real for
  symmetric \(C_\alpha\) and purely imaginary for antisymmetric \(C_\alpha\), the ring partition
  functions and the partition functions of the open chains with boundary operators \((b, b)\) at
  odd length and \((b, b^{\mathsf T})\) at even length satisfy the boundary moment inequalities at
  every \(\beta > 0\).
\end{theorem}
\begin{proof}
  \leanok
  Expanding the \(M\)-slice Trotter approximant of the open chain gives the alternating product
  \(c_L^{\mathsf T} \mathcal F^{\mathsf T} \mathcal F \mathcal F^{\mathsf T} \cdots c_R\).
  Reversing the imaginary-time sequences via the involution \(J\) satisfies
  \(\mathcal F^{\mathsf T} = J \mathcal F J\), so \(X = \mathcal F J\) is real symmetric and the
  \(n\)-site alternating product becomes \(J X^n J\) or \(J X^n\). Because the coefficient
  condition on \(b\) implies \(J c_L = \overline{c_L}\), both parities take the bilinear form
  \(\langle v, X^n v\rangle\) while the ring traces equal \(\Tr X^{2N}\). Passing to the
  \(M \to \infty\) limit via Lemma~\ref{lem:open-moment-bounds-limit} and the Lie--Trotter product
  formula completes the proof.
\end{proof}

\section{The open-chain spacetime-bootstrap criterion}

\begin{theorem}[Open-chain spacetime-bootstrap criterion]
  \label{thm:open-criterion}
  \lean{CertifiedGaps.OpenMomentWitness.move_one, CertifiedGaps.OpenMomentWitness.move_two,
    CertifiedGaps.OpenMomentWitness.oddChain_purity,
    CertifiedGaps.open_step, CertifiedGaps.ring_doubling_six, CertifiedGaps.gapGT_of_openSeeds,
    CertifiedGaps.RPRing.openGapGT_of_seeds}
  \leanok
  \uses{def:open-moment-bounds, lem:doubling-coefficient, lem:purity-squaring, lem:purity-to-gap,
    thm:open-moment-bounds}
  Assume the transfer-operator moment inequalities for the ring and the boundary moment
  inequalities at every \(\beta > 0\). If \(\Ps(120, 49) \ge 0.948\) and
  \(\Pt(240, 49) \ge 0.948\) for the ring, the boundary defect satisfies \(e_{142} \le 0.14\) at
  \(\beta = 49\), and the opposite-field open-chain thermal purity satisfies
  \(\Pt(240, 49) \ge 0.79\), then every odd open chain of \(n \ge 1921\) sites has a unique ground
  state and a spectral gap greater than \((\log 10)/392\).
\end{theorem}
\begin{proof}
  \leanok
  At \(\beta_i = 49 \cdot 2^i\), the ring defects satisfy \(\varepsilon_{i+1} = 6\varepsilon_i^2\)
  (Lemma~\ref{lem:doubling-coefficient} with \(C = 6\)) and control both the subleading spectral
  ratio and the bulk normalization ratio \(\mu_{i+1}/\mu_i^2\). Three recurrence steps advance the
  boundary defect and the open-chain thermal purity from scale \(i\) to scale \(i + 1\): the
  overlap recurrence, length doubling, and thermal squaring. Three explicit rational steps reach
  \(i = 3\), after which all defects lie below \(x_j = 10^{-1-2^j}\) with \(x_{j+1} = 10x_j^2\).
  At odd lengths the leading parity sign is \(s = +1\) because the partition function is positive,
  giving \(\Pt(n, \beta_i) \ge 1 - 4x\) and hence the uniform gap bound by
  Lemma~\ref{lem:purity-to-gap}.
\end{proof}

\section{The open spin-1 Heisenberg chain with boundary fields}

\begin{lemma}[Finite-size seed bounds for the boundary-field chain]
  \label{lem:open-seeds}
  \lean{CertifiedGaps.SpinOneHeisenberg.seed_ring_spatial,
    CertifiedGaps.SpinOneHeisenberg.seed_ring_thermal,
    CertifiedGaps.SpinOneHeisenberg.seed_boundary,
    CertifiedGaps.SpinOneHeisenberg.seed_open_thermal}
  \leanok
  \uses{def:heisenberg-spectrum, def:open-moment-bounds}
  At \(\beta = 49\), the periodic spin-1 Heisenberg ring satisfies \(\Ps(120) \ge 0.948\) and
  \(\Pt(240) \ge 0.948\), and the opposite-field open chain with \(h = 3/5\) satisfies
  \(e_{142} \le 0.14\) and \(\Pt(240) \ge 0.79\).
\end{lemma}
\begin{proof}
  These are the four finite-size numerical inputs of~\cite{OpenAIBoundary}, which enter the
  formalization as explicit hypotheses.
\end{proof}

\begin{theorem}[Haldane gap of the open spin-1 chain with boundary fields]
  \label{thm:boundary-field-gap}
  \lean{CertifiedGaps.SpinOneHeisenberg.endFieldHamiltonian,
    CertifiedGaps.SpinOneHeisenberg.openHamiltonian_odd,
    CertifiedGaps.SpinOneHeisenberg.openHamiltonian_even,
    CertifiedGaps.SpinOneHeisenberg.gapGT_of_boundaryFieldInputs,
    CertifiedGaps.SpinOneHeisenberg.spin_one_boundary_field_gap}
  \leanok
  \uses{thm:open-criterion, lem:open-seeds, thm:heisenberg-moments}
  For every odd \(n \ge 1921\), the open spin-1 Heisenberg chain with equal boundary fields
  \(h = 3/5\) has a unique ground state, and every excited eigenvalue lies more than
  \((\log 10)/392\) above the ground energy.
\end{theorem}
\begin{proof}
  \leanok
  In the Cartesian basis, taking the imaginary coefficient \(w_z = -ih\) on the antisymmetric
  generator \(G_z = i S^z\) produces the boundary operator \(b = -h S^z\), whose transpose is
  \(b^{\mathsf T} = +h S^z\); the open chain of Theorem~\ref{thm:open-moment-bounds} therefore has
  equal boundary fields at odd length and opposite boundary fields at even length. Applying
  Theorem~\ref{thm:open-criterion} with the inputs of Lemma~\ref{lem:open-seeds} gives the bound.
\end{proof}
